% Includable research supplement. Theorem environments and bibliography are
% provided by paper.tex. All variable counts include the input unless stated.

\section{Stronger obstructions for small polynomial presentations}
\label{p2lb:scope}

The unrestricted bounds established in the first report are two existential
witnesses and total degree four.  The arguments below give stronger bounds
for important classes of presentations.  They do not prove that every
definition needs more than two witnesses.  In particular, the general case
of a sign-changing quartic in three total variables remains unresolved here.

For a number field $K$, write $\mathcal O_K$ for its ring of integers.
If $V\subseteq\mathbb A^n_K$ is an algebraic set with first coordinate $T$,
put
\[
 \pi(V(K))=\{t\in K:\text{some }(t,y_1,\ldots,y_{n-1})\in V(K)\}.
\]
Thus $n$ counts all coordinates, and $n-1$ counts existential witnesses.
All assertions in this supplement are independent of the proposed
arithmetic proof in the repository.

\subsection{Curves and rational parametrizations}

We use the following consequences of Siegel's theorem and the topology
of a number field.  They apply to open subsets, which is important when
using birational maps.

\begin{lemma}\label{p2lb:curves}
Let $V$ be an algebraic set over $K$ of dimension at most one, and let
$f:V\to\mathbb A^1_K$ be a regular function.  If
$f(V(K))\subseteq\mathcal O_K$, then $f(V(K))$ is finite.
The same conclusion holds for a rational function restricted to any
Zariski-open subset on which it is defined.
\end{lemma}
\begin{proof}
There are finitely many irreducible components.  A zero-dimensional
component contributes finitely many points, and a component on which
$f$ is constant contributes at most one value.  On an irreducible curve
which is not geometrically irreducible, every $K$-point lies on the
intersection of two distinct conjugate geometric components.  There
are only finitely many such points.

For a geometrically integral curve, normalization and projective
completion discard or add only finitely many points.  Let $C$ be the
resulting smooth projective curve and let $f\in K(C)$ be nonconstant.
If $C$ has positive genus, Siegel's theorem in its rational-function
form gives only finitely many $P\in C(K)$ for which $f(P)$ is defined
and integral.  One obtains this form from the usual theorem by removing
the poles of $f$ and enlarging a finite set of primes so that the
coordinate functions of that affine curve are integral over
$\mathcal O_{K,S}[f]$; see \cite[Chapter~8]{lbLang1983}.

If $C$ has genus zero and has a $K$-point, it is $K$-isomorphic to
$\mathbb P^1$.  A nonconstant rational function on a nonempty open
subset of $\mathbb P^1(K)$ cannot take only integral values.  Indeed,
embed $K$ diagonally in
$K_\infty=\mathbb R^{r_1}\times\mathbb C^{r_2}$.  The image of $K$
is dense, while that of $\mathcal O_K$ is a discrete lattice.  Continuity
would force the rational function to be constant on a small connected
open ball, and hence identically constant.  If $C(K)$ is empty, only
the previously discarded singular points remain.  These cases prove
the assertion.
\end{proof}

\begin{lemma}\label{p2lb:unirational}
Let $X/K$ be an integral unirational variety, let $U\subseteq X$ be
nonempty and open, and let $f\in K(X)$ be nonconstant and defined on $U$.
Then $f(U(K))$ is not contained in $\mathcal O_K$.
\end{lemma}
\begin{proof}
Compose $f$ with a dominant rational map from projective space to $X$.
After restricting to an affine open set, the composition is a
nonconstant rational function in finitely many variables.  The density
and discreteness argument in the preceding proof, now on
$K_\infty^m$, gives a contradiction if all its values are integral.
\end{proof}

\begin{lemma}\label{p2lb:ruled}
Let $C/K$ be a geometrically integral curve, let $m\geq1$, and let
$X$ be birational over $K$ to $C\times\mathbb P^m$.
Suppose $U\subseteq X$ is nonempty and open and $U(K)$ is Zariski
dense.  A nonconstant $f\in K(X)$, defined on $U$, cannot satisfy
$f(U(K))\subseteq\mathcal O_K$.
\end{lemma}
\begin{proof}
Restrict the birational correspondence to isomorphic nonempty opens
and replace $C$ by its smooth projective model.  We may work on an
open subset of $C\times\mathbb A^m$.  If $f$ is nonconstant in the
affine coordinates over $K(C)$, then outside finitely many points of
$C$ its restriction to a fiber is nonconstant.  Outside another finite
set that fiber meets the chosen open subset in a nonempty open subset
of $\mathbb A^m$.  Density of $U(K)$ supplies a rational base point
outside these exceptional sets.  The restriction to this fiber violates
Lemma~\ref{p2lb:unirational}.

Otherwise $f$ belongs to $K(C)$ and is nonconstant there.  For every
$c\in C(K)$ outside a finite set, the open fiber has a $K$-point,
so $f(c)$ is integral.  Lemma~\ref{p2lb:curves} gives only finitely
many such values, hence only finitely many such $c$, because a
nonconstant map of projective curves has finite fibers.  This
contradicts the density of $U(K)$ in $C\times\mathbb A^m$.
\end{proof}

\begin{remark}
The preceding lemma concerns a variety birational over $K$ to a
product with projective space.  Geometric ruledness after extending
scalars is not sufficient for its proof.  In particular, it must not
be applied without further argument to an arbitrary conic bundle
over $\mathbb P^1_K$ which has no rational section.
\end{remark}

\begin{lemma}\label{p2lb:dense-component}
Let $V/K$ be an algebraic set and $f:V\to\mathbb A^1_K$ a regular
function.  If $f(V(K))$ is infinite and contained in $\mathcal O_K$,
the Zariski closure of $V(K)$ has a geometrically integral component
$X$ of dimension at least two such that $X(K)$ is Zariski dense
and $f|_X$ is nonconstant.
\end{lemma}
\begin{proof}
The closure has finitely many irreducible components, each with dense
$K$-points on the open set obtained by deleting the other components.
If such a component were geometrically reducible, its rational points
would lie in the proper closed intersections of its conjugate
geometric components, a contradiction.  Components on which $f$ is
constant contribute finitely many values, and components of dimension
at most one contribute finitely many by Lemma~\ref{p2lb:curves}.
Thus at least one component has the required properties.
\end{proof}

\subsection{Two quadrics in projective four-space}

The next argument uses only the elementary degree bound
$\deg X\geq1+\operatorname{codim}(X)$ for a nondegenerate projective
integral variety, B\'ezout's theorem, and the following theorem of
Koll\'ar: a geometrically integral cubic hypersurface of dimension at
least two over a perfect field is unirational if it has a rational
point and is not a cone \cite[Theorem~2]{lbKollar2002}.
The elementary degree bound is recalled in
\cite[Introduction]{p2lbEGHP2006}.  We give the projection argument
explicitly in order to include reducible and nonreduced intersections.

\begin{proposition}\label{p2lb:two-quadrics-surface}
Let $Q_1,Q_2\in K[X_0,\ldots,X_4]$ be coprime homogeneous
quadratics, and let $X$ be a geometrically integral surface component
of $V(Q_1,Q_2)$.  Suppose $X(K)$ is Zariski dense.
Then $X$ is either unirational over $K$ or birational over $K$ to
$C\times\mathbb P^1$ for a geometrically integral curve $C/K$.
\end{proposition}
\begin{proof}
The complete intersection $W=V(Q_1,Q_2)\subseteq\mathbb P^4$
has degree four.  More precisely, if $m_X$ is its multiplicity along
$X$, then $m_X\deg X\leq4$.  This follows either from B\'ezout's
theorem with multiplicities or from the complete-intersection Hilbert
series $(1-t^2)^2/(1-t)^5$ and additivity of degree.

If $\deg X\leq2$, the elementary degree bound shows that its linear
span has dimension at most three.  It is therefore a plane or an
integral quadric surface in its span.  Its linear span is defined over
$K$, and density supplies a smooth $K$-point in the quadric case.
Projection from that point makes the quadric rational; a plane is
rational as well.

Suppose $\deg X\geq3$.  The degree inequality forces $m_X=1$.
Since the characteristic is zero, a nonempty open subset of $X$ is
smooth in the complete intersection $W$.  Choose a $K$-point there,
using density, and change coordinates so that it is
$p=(1:0:0:0:0)$.  The quadratics have the form
\[
 Q_1=X_0L_1+q_1,\qquad Q_2=X_0L_2+q_2,
\]
where $L_i$ are linear and $q_i$ are quadratic in $X_1,\ldots,X_4$.
Smoothness of $W$ at $p$ means that $L_1,L_2$ are independent.
The cubic
\[
 H=L_1q_2-L_2q_1
\]
is not identically zero.  Otherwise coprimality of $L_1,L_2$ gives
$q_1=L_1L_3$ and $q_2=L_2L_3$ for a linear form $L_3$, contradicting
coprimality of $Q_1,Q_2$.

The surface $X$ cannot be contained in the plane $L_1=L_2=0$,
because its degree is at least three.  Interchange the two quadratics
if necessary so that $L_1$ does not vanish identically on $X$.
Projection from $p$ then has the rational inverse
\[
 X_0=-q_1/L_1
\]
on the corresponding open set.  Thus $X$ is birational to a
geometrically integral surface component $Y$ of the cubic
$V(H)\subseteq\mathbb P^3$.  The rational points of $Y$ are dense.

If $Y$ is a plane or a quadric, it is rational as above.  If it is
a cubic which is not a cone, Koll\'ar's theorem makes it unirational.
Finally, a geometrically integral cubic surface which is a cone has
a single vertex point defined over $K$.  Indeed the vertex directions
form the kernel of a linear map over $K$ given by the directional
derivatives of its cubic equation; more than one independent vertex
direction would leave a binary cubic, which is geometrically reducible.
The complement of the vertex is birational over $K$ to the product
of its plane cubic base with $\mathbb P^1$.  This proves the result.
\end{proof}

\begin{theorem}\label{p2lb:quadratic-system-four}
Let $q_1,\ldots,q_s\in K[T,Y_1,\ldots,Y_{n-1}]$ have total
degree at most two, where $n\leq4$.  If the rational projection
\[
 S=\{t\in K:\exists\mathbf y\in K^{n-1},\quad
                   q_1(t,\mathbf y)=\cdots=q_s(t,\mathbf y)=0\}
\]
is contained in $\mathcal O_K$, then $S$ is finite.
\end{theorem}
\begin{proof}
Adding unused coordinates reduces the assertion to $n=4$.
Zero residuals can be removed; if all residuals vanish identically,
then $S=K$, contrary to the hypothesis.  A nonzero constant residual
gives the empty set.  Let $V$ be their common zero set.

Suppose $S$ is infinite, and choose $X\subseteq V$ as in
Lemma~\ref{p2lb:dense-component}.  Any common polynomial factor
$g$ of the residuals has degree at most two.  The entire hypersurface
$g=0$ lies in $V$, and its input projection is therefore contained
in $\mathcal O_K$.  The quadratic obstruction proved in the first
report says this projection has at most two elements.  Consequently
$X$ is not contained in $g=0$.  Dividing the residuals by their common
greatest divisor does not remove $X$ and leaves quadratics with
greatest common divisor one.

If their $K$-linear span contains a nonzero affine linear polynomial,
restrict to its affine hyperplane, isomorphic to $\mathbb A^3_K$.
Not all restricted residuals vanish identically, since that would
give the original quadratics a common linear factor.  The resulting
surface $X$ is a hypersurface component of a nonzero quadratic in
three variables, hence is a plane or a quadric.  Density gives a
rational parametrization, contradicting Lemma~\ref{p2lb:unirational}.

Otherwise choose two coprime quadratics in the linear span of the
residuals.  To see that this is possible, first choose a nonzero
quadratic $q$ and then a linear combination not divisible by any
irreducible factor of $q$; finitely many proper linear subspaces
cannot cover a vector space over the infinite field $K$.
Their affine common zero set has dimension two.  Thus $X$ is a
surface component.  Their degree-two homogenizations remain coprime,
and the projective closure of $X$ is a component of their intersection
in $\mathbb P^4$.

Proposition~\ref{p2lb:two-quadrics-surface} makes this component
unirational or birational to a product of a curve and $\mathbb P^1$.
The input is a nonconstant rational function and is integral on
the rational points of its affine open subset.  This contradicts
Lemma~\ref{p2lb:unirational} or Lemma~\ref{p2lb:ruled}.
\end{proof}

\begin{corollary}\label{p2lb:sos-five}
Suppose
\[
 P(T,\mathbf Y)=\sum_{j=1}^s c_jq_j(T,\mathbf Y)^2,
 \qquad c_j\in\mathbb Q_{>0},\quad\deg q_j\leq2.
\]
If the rational existential projection of $P=0$ is $\mathbb Z$,
then $P$ uses at least five total variables, hence at least four
existential witnesses.
The same lower bound holds over a number field for any quartic
obtained by applying an anisotropic quadratic form to quadratic
residuals and defining its ring of integers.
\end{corollary}
\begin{proof}
In the rational case, the sum is zero exactly when all $q_j$ vanish.
For an anisotropic combining form the same equivalence holds by
definition of anisotropy.  Apply Theorem~\ref{p2lb:quadratic-system-four}.
\end{proof}

This improves the unrestricted three-total-variable bound for the
sum-of-squares quartic architecture used by the released polynomial
and all the expanded quartic variants in the first report.  It does
not rule out a sign-changing quartic in three or four total variables.

\subsection{Nonnegative polynomials in three variables}

\begin{proposition}\label{p2lb:nonnegative-squarefree}
Let $P\in\mathbb Q[T,Y,Z]$ be squarefree and nonnegative on
$\mathbb R^3$.  If its rational input projection is contained in
$\mathbb Z$, that projection is finite.
\end{proposition}
\begin{proof}
Every real zero of a differentiable nonnegative function is a local
minimum, so all three partial derivatives of $P$ vanish there.
The singular locus of a squarefree hypersurface in characteristic
zero has codimension at least one in that hypersurface.  In this
case it has dimension at most one.  All rational zeros lie in this
singular locus.  Apply Lemma~\ref{p2lb:curves} to its input function.
\end{proof}

\begin{theorem}\label{p2lb:nonnegative-degree-eight}
Let $P\in\mathbb Q[T,Y,Z]$ be nonnegative on $\mathbb R^3$.
If its rational input projection is an infinite subset of $\mathbb Z$,
then $P$ has an irreducible factor of degree at least four and
multiplicity at least two.  In particular $\deg P\geq8$.
\end{theorem}
\begin{proof}
Choose a component $X$ of the closure of its rational zeros as in
Lemma~\ref{p2lb:dense-component}.  It has dimension two, so its prime
ideal is generated by an irreducible factor $g$ of $P$.  Write
$P=g^m h$ with $g\nmid h$.  Density supplies a rational point on
$X$ which is smooth on $g=0$ and does not lie on $h=0$.
If $m=1$, then $\nabla P=h\nabla g$ is nonzero there, contrary
to nonnegativity of $P$.  Thus $m\geq2$.

The projection of the whole hypersurface $g=0$ is contained in
the projection of $P=0$, and is infinite because the input is
nonconstant on $X$ and $X(\mathbb Q)$ is dense.  The unconditional
degree-three obstruction of the first report gives $\deg g\geq4$.
Therefore $\deg P\geq m\deg g\geq8$.
\end{proof}

Nonnegativity only on rational inputs gives the same result, because
$\mathbb Q^3$ is dense in $\mathbb R^3$ and $P$ is continuous.
The degree-eight conclusion is restricted to nonnegative polynomials;
it is not an unconditional degree-eight lower bound for two witnesses.

\subsection{Further necessary conditions}

\begin{proposition}\label{p2lb:linear-coordinate}
If a polynomial $P(T,Y_1,\ldots,Y_r)\in K[T,\mathbf Y]$ defines
$\mathcal O_K$ in $K$, its degree in each witness which occurs is
at least two.  Its degree in $T$ is also at least two.
\end{proposition}
\begin{proof}
Suppose $P=AY_r+B$, where $A\ne0$ is independent of $Y_r$.
For all but finitely many $t\in K$, the specialization
$A(t,Y_1,\ldots,Y_{r-1})$ is a nonzero polynomial.  Since $K$ is
infinite, choose the other witnesses so that it does not vanish and
then solve uniquely for $Y_r$.  The projection is cofinite in $K$,
which is incompatible with $\mathcal O_K$.

If $P=AT+B$ with $A\ne0$ independent of $T$, its projection includes
the values of $-B/A$ on the nonempty open set $A\ne0$.  Integrality
of all these values forces this rational function to be constant,
by the density and discreteness argument.  Thus $B=-cA$ for some
$c\in K$, and $P=A(T-c)$.  If $A$ has a $K$-zero, every input is
accepted; otherwise only $c$ is accepted.  Neither is $\mathcal O_K$.
If $P$ is independent of $T$, its projection is empty or all of $K$.
\end{proof}

\begin{proposition}\label{p2lb:quartic-surface-necessary}
Suppose a quartic $P\in\mathbb Q[T,Y,Z]$ defines $\mathbb Z$ in
$\mathbb Q$.  Then $P$ is irreducible and geometrically irreducible,
its rational points are Zariski dense on its affine surface, and
$P$ takes both positive and negative values on $\mathbb R^3$.
The input function is nonconstant on that surface.  Every curve
birational over $\mathbb Q$ to $\mathbb P^1$ on which the input is defined is vertical
(the input is constant), and the same holds for every geometrically
integral genus-one curve on the surface with infinitely many rational
points.
\end{proposition}
\begin{proof}
If $P$ factors nontrivially, every factor has degree at most three.
Each factor's rational projection is contained in $\mathbb Z$ and
is finite by the earlier degree-three obstruction.  Their finite
union cannot be $\mathbb Z$.  If the irreducible $P$ is geometrically
reducible, its rational points lie in intersections of distinct
conjugate components, of dimension at most one; Lemma~\ref{p2lb:curves}
again gives a finite projection.  The same lemma shows that a proper
Zariski closure of the rational points would give a finite projection.
Thus the surface is geometrically integral and its rational points
are dense.  The input is nonconstant because its image is infinite.

If $P$ had only one sign, apply
Theorem~\ref{p2lb:nonnegative-degree-eight} to $P$ or $-P$.
Finally, a nonvertical rational curve with a rational point would
violate Lemma~\ref{p2lb:unirational}.  On a genus-one curve with
infinitely many rational points a nonconstant input function would
take only finitely many integral values at rational points by
Siegel's theorem, whereas every rational point where the input is
defined would have to give such a value.  This is impossible.
\end{proof}

\subsection{Why the unrestricted two-witness case remains difficult}

A smooth quartic surface in $\mathbb P^3$ has trivial canonical
bundle and is a K3 surface.  It is not covered by the unirational
cubic theorem.  Zariski density of rational points, even when known,
does not itself imply density in the real topology or the existence
of a nonintegral value of a specified rational function.  Likewise,
potential density after a finite field extension cannot be substituted
for a statement about rational points over $\mathbb Q$.
For examples and the distinction between density over the ground field
and potential density, see \cite{p2lbHarrisTschinkel2000}.

The same limitation appears naturally in the improved bound for
quadratic systems.  A smooth complete intersection of three quadrics
in $\mathbb P^5$ is a K3 surface: adjunction gives
$K_X=\mathcal O_X(2+2+2-6)=\mathcal O_X$, and the
complete-intersection cohomology gives $H^1(X,\mathcal O_X)=0$.
Thus five affine coordinates already permit a class of surfaces not
handled by the two-quadric argument.  This observation is a limit of
the proof, not an existence claim for a five-variable definition.

\begin{proposition}[Conditional obstruction for general-type surfaces]
\label{p2lb:lang-surface}
Assume the weak Bombieri--Lang conjecture for surfaces over $\mathbb Q$:
the rational points of a smooth projective surface of general type are
not Zariski dense.  Then no polynomial in three total variables whose
geometrically integral projective surface has a smooth projective
model of general type can have an infinite rational input projection
contained in $\mathbb Z$.
In particular, under this conjecture a polynomial in three variables
with smooth projective closure of degree at least five cannot define
$\mathbb Z$.
\end{proposition}
\begin{proof}
Lemma~\ref{p2lb:curves} forces Zariski density of rational points
on a surface with such a projection.  Passing to isomorphic dense
opens of a smooth projective model preserves density, contradicting
the conjecture.  A smooth degree-$d$ surface in $\mathbb P^3$ has
canonical bundle $\mathcal O_X(d-4)$, which is ample when $d\geq5$.
\end{proof}

Mazur's conjecture about finitely many real connected components of
the closure of rational points would rule out every Diophantine
definition of $\mathbb Z$ in $\mathbb Q$, in any number of variables.
It therefore cannot supply a conditional numerical improvement
compatible with simultaneously assuming the proposed arithmetic
theorem.  Proposition~\ref{p2lb:lang-surface}, by contrast, is a
restriction on one class of possible two-witness surfaces.
